% SPDX-License-Identifier: CC-BY-4.0
% P3 common-mode integrity manuscript. Target: NAVIGATION (ION) / IEEE Trans. AES.
% Build: tectonic main.tex   (or pdflatex + bibtex + pdflatex x2)
\documentclass[11pt]{article}

\usepackage[utf8]{inputenc}
\usepackage[T1]{fontenc}
\usepackage{lmodern}
\usepackage[margin=1in]{geometry}
\setlength{\emergencystretch}{3em}
\usepackage{amsmath}
\usepackage{amssymb}
\usepackage{amsthm}
\usepackage{bm}
\usepackage{graphicx}
\usepackage{booktabs}
\usepackage{array}
\usepackage{siunitx}
\usepackage{xcolor}
\usepackage[colorlinks=true,linkcolor=black,citecolor=black,urlcolor=blue]{hyperref}
\usepackage[numbers,sort&compress]{natbib}
\usepackage{orcidlink}

\graphicspath{{./}{figures/}}
\setlength{\abovecaptionskip}{9pt}
\setlength{\belowcaptionskip}{6pt}

% Theorem environments
\newtheorem{theorem}{Theorem}
\newtheorem{lemma}{Lemma}
\newtheorem{proposition}{Proposition}
\newtheorem{corollary}{Corollary}
\theoremstyle{definition}
\newtheorem{definition}{Definition}
\newtheorem{remark}{Remark}

% Notation macros
\newcommand{\code}[1]{\texttt{#1}}
\newcommand{\R}{\mathbb{R}}
\newcommand{\Gmat}{\bm{G}}                      % geometry (design) matrix
\newcommand{\yv}{\bm{y}}                        % measurement-residual vector
\newcommand{\xv}{\bm{x}}                        % state (position + clock)
\newcommand{\Smat}{\bm{S}}                      % parity matrix
\newcommand{\Pperp}{\bm{P}^{\!\perp}}          % parity (residual) projector
\newcommand{\Ppar}{\bm{P}_{\mathcal{R}}}       % range projector
\newcommand{\los}{\bm{e}}                       % line-of-sight unit vector
\newcommand{\sv}{\bm{s}}                        % satellite position
\newcommand{\uv}{\bm{u}}                        % user position
\newcommand{\ds}{\Delta\bm{s}}                  % common frame/ephemeris shift
\newcommand{\rng}{\operatorname{range}}
\newcommand{\Sig}{\bm{\Sigma}}                  % covariance
\DeclareMathOperator{\rank}{rank}

% Honesty tags for proven / validated / representative-modelled quantities
\newcommand{\Proven}{\textsc{Proven}}
\newcommand{\Modelled}{\textsc{Modelled}}
\newcommand{\Validated}{\textsc{Validated}}
\newcommand{\modtag}{{\footnotesize\textsf{[M]}}}

\title{\textbf{Common-mode integrity for multi-provider lunar PNT: the
reference-frame and ephemeris error class RAIM/ARAIM cannot see, and a
cross-provider protection level}}

\author{Chakshu Baweja\,\orcidlink{0009-0008-2098-0751}\\
\small Ashforde O\"U, Tallinn, Estonia \\
\small \texttt{contact@ashforde.org}}

\date{\today}

\begin{document}
\maketitle

\begin{abstract}
A multi-provider lunar positioning, navigation, and timing (PNT) architecture---the
declared vision of LunaNet, ESA Moonlight/LCNS, and their national counterparts---will
rely on receiver autonomous integrity monitoring (RAIM), and its advanced variant
ARAIM, to bound position error in real time. We study one error class that the
\emph{snapshot residual test} at the heart of RAIM/ARAIM cannot see, and we are
careful throughout to separate what is classical, what we prove, what we validate on
real data, and what we model. First, a classical fact we cite rather than claim: the
parity-space formulation detects a fault only through the parity subspace (the
orthogonal complement of $\rng(\Gmat)$), so any measurement error lying in
$\rng(\Gmat)$ produces a zero parity vector and is undetectable
\citep{sturza1988,brown1992,parkinson1988}. Second, a lemma we prove: to first order
the \emph{entire} rigid-body transformation of the satellite constellation---three
translations and three rotations, and independent of the user's own position---lies in
$\rng(\Gmat)$ and is therefore fully blind; only non-rigid deformation is detectable,
and the residual of a rigid rotation is second order in the angle
($\sim\!10^{-9}\,$m for realistic nanoradian frame errors). This is exactly the
geodetic datum defect / gauge freedom \citep{baarda1973,teunissen2000}, so---connecting
to our companion datum-identifiability paper \citep{baweja2026datum}---the P1
datum-defect subspace is a subset of the RAIM blind subspace. We claim novelty for
none of these three (subspace blindness, the datum-defect$\subset$blind relation, and
bounding an undetectable error by a protection level are all prior art). Third, the one
externally \Validated{} row: on the real geocentric-Moon disagreement between the
independent authoritative ephemerides DE440 (JPL), INPOP21a (IMCCE), and EPM2021
(IAA~RAS), the engine's parity/least-squares split reproduces an independent NumPy
computation to relative and absolute error below $10^{-3}$
($732$ records $=2$ provider pairs $\times366$ epochs). We are scrupulous about what
realness contributes: the blindness itself (blind fraction $=1$, parity residual
$\sim\!10^{-15}\,$m) is an \emph{analytic identity} of the common-shift construction
$\delta\yv=\Gmat[-\ds;0]$, not a measurement---the $\sim\!10^{-15}\,$m is
floating-point zero, and the data set only the $\approx\!2\,$m \emph{magnitude} of the
absorbed user-position error. Fourth, and \Modelled{}: mapped through a modelled lunar
constellation this $\approx\!2\,$m inter-provider common-mode motivates a
\emph{cross-provider protection level} (CMPL). Here we state the load-bearing claim
correctly, because its naive form is false: ARAIM is \emph{not} ``provably blind'' to
common-mode at the protection level---the ARAIM protection level already bounds
residual-invisible error through its fault-free hypothesis, via the $\sigma_{\mathrm
{URA}}$ overbound and the nominal bias $b_{\mathrm{nom}}$
\citep{blanch2015,wgc2016araim}. The genuinely uncovered term is the multi-provider
case: a frame/ephemeris common-mode between \emph{independent} lunar providers lies
outside any single provider's per-provider URA/ISM budget. The CMPL restores a valid
bound for exactly that case, and the conservative envelope
$\mathrm{HPL}_{\mathrm{total}}=\mathrm{HPL}_{\mathrm{ARAIM}}+\mathrm{CMPL}$ is a
triangle-inequality sum (a root-sum-square would under-bound and overclaim integrity),
valid only where the common-mode is not already absorbed into $\sigma_{\mathrm{URA}}$,
else it double-counts. The classical property and the rigid-body lemma are proven; the
parity/least-squares split of the real ephemeris disagreement is \Validated{}; the
constellation, the common-mode covariance, and the whole CMPL envelope are explicitly
\Modelled{} (the ``providers'' are ephemerides, a representative analog---no two
lunar-PNT providers fly), and every number is reproducible from a single command.

\end{abstract}

\noindent\textbf{Keywords:} lunar PNT; integrity; RAIM; ARAIM; protection level;
common-mode error; parity space; reference frame; planetary ephemerides; DE440;
INPOP; EPM; LunaNet; Moonlight

\vspace{0.5em}
\section{Introduction}
\label{sec:intro}

The next decade will place several independent positioning, navigation, and timing
(PNT) services around the Moon. NASA's LunaNet and its Lunar Communications Relay and
Navigation Systems, ESA's Moonlight programme and its Lunar Communications and
Navigation Service (LCNS), and Japan's Lunar Navigation Satellite System are each
designed to broadcast navigation signals to surface and orbital users
\citep{israel2020lunanet,fienga2024moonlight}. Safety-relevant lunar operations---%
descent and landing, surface mobility, rendezvous---will need not only a position but
a real-time \emph{bound} on that position's error: an integrity monitor. On Earth this
is the role of receiver autonomous integrity monitoring (RAIM), and of its advanced
multi-constellation successor ARAIM \citep{blanch2015,wgc2016araim}, which combine a
snapshot \emph{residual test} that detects faulty measurements with a
\emph{protection level} that bounds the position error to a stated integrity risk.
Both are natural candidates for lunar PNT, and both are already assumed in
lunar-navigation performance studies.

This paper isolates one error class that the snapshot residual test at the heart of
RAIM/ARAIM structurally cannot detect, quantifies it on real data, and shows what a
multi-provider lunar architecture must add to bound it. We are deliberate, throughout,
about the status of each claim---classical, proven, validated, or modelled---because
the naive version of the headline is wrong and a careful reviewer will (rightly)
reject it.

\paragraph{The classical fact, cited not claimed.}
Snapshot RAIM/ARAIM forms a linearized measurement model $\yv=\Gmat\xv+\bm\varepsilon$
and detects a fault through the \emph{parity space}: the orthogonal complement of the
column space $\rng(\Gmat)$ of the geometry matrix. The fault-detection statistic is the
norm of the parity vector, and a measurement error $\delta\yv$ enters it only through
its parity projection $\Pperp\delta\yv$. Consequently \emph{any} error lying in
$\rng(\Gmat)$ has a zero parity vector and is invisible to the test---it is absorbed,
without residual, into the estimated state. This is textbook: it is the parity-space
formulation of \citet{sturza1988}, the RAIM equivalence of \citet{brown1992}, and the
pseudorange-residual monitor of \citet{parkinson1988}, and it is even noted in our own
engine's parity-RAIM detector. We claim no novelty for it.

\paragraph{What we prove.}
The classical statement is about an abstract subspace; our first contribution makes it
concrete for the constellation. We prove (Lemma~\ref{lem:rigid}) that to first order
the \emph{entire} rigid-body transformation of the satellite constellation---three
translations and three rotations, applied in common to every satellite, and
\emph{independent of the user's own position}---lies in $\rng(\Gmat)$ and is therefore
fully blind. Only \emph{non-rigid} deformation of the constellation is detectable, and
the residual of a rigid rotation is second order in the angle
($\propto|\bm\theta|^2$, on the order of $10^{-9}\,$m for the nanoradian-scale frame
errors realistic between lunar reference realizations). This rigid-body subspace is
exactly the geodetic \emph{datum defect}, or gauge freedom, of free-network adjustment
\citep{baarda1973,teunissen2000}. It follows---and this is a synthesis, not a
discovery---that the single-provider datum-defect subspace classified in our companion
paper \citep{baweja2026datum} is a \emph{subset} of the RAIM blind subspace: an error
no provider's own network can fix is also an error that provider's integrity monitor
cannot see.

\paragraph{What we validate on real data.}
The ``multiple providers'' problem exists already, one level down, in a real and
measurable form. Three independent institutions maintain authoritative lunar and
planetary ephemerides from independent data reductions and dynamical models: DE440 from
JPL/NASA \citep{park2021de440}, INPOP21a from IMCCE \citep{fienga2021inpop21a}, and
EPM2021 from IAA~RAS \citep{pitjeva2021epm2021}. Their pairwise geocentric-Moon
disagreement---the same real quantity our companion cross-provider interoperability
paper decomposes \citep{baweja2026interop}---is the closest existing analogue of an
inter-provider frame/ephemeris common-mode, measurable today with no flight data. We
take that disagreement as a common shift $\ds$ of the constellation, form the
measurement error $\delta\yv=\Gmat[-\ds;0]$, and split it into its range and parity
parts. Our one externally \Validated{} row is that the engine's Rust parity/%
least-squares split reproduces an independent NumPy computation to relative and
absolute error below $10^{-3}$, over $732$ records ($2$ provider pairs $\times366$
epochs). We are scrupulous about what realness buys: the \emph{blindness} itself---%
blind fraction $=1$, parity residual $\sim\!10^{-15}\,$m---is an analytic identity of
the common-shift construction, guaranteed by $\delta\yv\in\rng(\Gmat)$ and \emph{not}
produced by the data; the $\sim\!10^{-15}\,$m is floating-point zero, not a measured
residual. What the data contribute is only the $\approx\!2\,$m \emph{magnitude} of the
resulting user-position error. The \Validated{} content is therefore precise and
narrow: the Rust linear algebra matches NumPy at real magnitudes.

\paragraph{What we model, and the load-bearing claim stated correctly.}
The lunar constellation geometry, the common-mode covariance, the cross-provider
protection-level (CMPL) magnitudes, a representative ARAIM horizontal protection level
(HPL), and the whole integrity envelope are \Modelled{} (the ``providers'' here map
real ephemerides through a modelled constellation---no two lunar-PNT providers fly).
The load-bearing claim must be stated in its correct form, because the naive form is
false. ARAIM is \emph{not} ``provably blind'' to common-mode at the protection level.
The RAIM/ARAIM \emph{residual test} is provably blind to $\rng(\Gmat)$; but the ARAIM
\emph{protection level} already bounds undetectable error through its fault-free
(nominal) hypothesis---the user range accuracy overbound $\sigma_{\mathrm{URA}}$ plus
the nominal bias $b_{\mathrm{nom}}$ of the multiple-hypothesis solution-separation
(MHSS) algorithm \citep{blanch2015,wgc2016araim}. Residual-invisible errors are still
budgeted, up to what the integrity support message (ISM) declares. The genuinely
uncovered term, and this paper's contribution, is the \emph{multi-provider} case: a
cross-provider frame/ephemeris common-mode between \emph{independent} lunar providers
lies outside \emph{any single provider's} per-provider URA/ISM budget. The CMPL
restores a valid bound for exactly that case, and the total envelope
$\mathrm{HPL}_{\mathrm{total}}=\mathrm{HPL}_{\mathrm{ARAIM}}+\mathrm{CMPL}$ is a
\emph{conservative} triangle-inequality sum---a root-sum-square would under-bound and
overclaim integrity, since the two contributions can align in the position domain---%
valid only where the common-mode is not already absorbed into $\sigma_{\mathrm{URA}}$
(otherwise it double-counts). We state this conditionality explicitly wherever the
envelope appears.

\paragraph{Honest-novelty statement.}
To fix the boundary once, in the words we use again in the discussion: \emph{the
blindness of snapshot RAIM/ARAIM to errors in $\rng(\Gmat)$ is a classical property of
the parity-space formulation and is not claimed as new. Our contribution is (i) the
real-data quantification that the actual median $2.0$--$2.4\,$m inter-ephemeris
disagreement between independent authoritative lunar ephemerides (DE440 vs.\ INPOP21a
and EPM2021), projected through a lunar constellation, is absorbed almost entirely as
user position error (blind fraction $\approx1$, parity residual $\sim\!10^{-15}\,$m)
and is therefore structurally invisible to the RAIM/ARAIM residual test; (ii) the
observation that in a multi-provider lunar PNT architecture this inter-provider
common-mode lies outside any single provider's per-provider URA/ISM budget; and (iii)
the CMPL, an additive position-domain term that restores a valid bound for exactly that
uncovered cross-provider case. We do not claim novelty for the subspace blindness, for
the datum-defect$\,\subset\,$blind relation, or for bounding an undetectable error by a
protection level---the last being the mechanism of ARAIM's own fault-free/nominal-bias
term.}

\paragraph{Honesty scope.}
We separate four kinds of statement, and label them at every occurrence. The classical
parity-space blindness (Proposition~\ref{prop:blind}) is \emph{prior art}, cited. The
rigid-body lemma and the datum-defect$\,\subset\,$blind relation
(Lemma~\ref{lem:rigid}, Corollary~\ref{cor:datum}) are \Proven{} algebra. The
parity/least-squares split of the real DE440/INPOP21a/EPM2021 disagreement
(Section~\ref{sec:anchor}) is \Validated{} against an independent numerical oracle, with
the explicit caveat that only its \emph{magnitude} is empirical, its blindness being an
analytic identity. Everything at the system level---the constellation, the common-mode
covariance $\Sig_{\mathrm{cm}}$, the CMPL, the representative HPL, and the total
envelope (Section~\ref{sec:cmpl})---is \Modelled{} and tagged \modtag{}, because the
providers are ephemerides. Section~\ref{sec:validation} makes the firewall precise, and
every number is regenerated from the committed engine by a single command
(Section~\ref{sec:availability}).

\section{The snapshot measurement model and the parity split}
\label{sec:model}

We work with the standard linearized single-epoch (snapshot) pseudorange model used by
RAIM and ARAIM \citep{parkinson1988,brown1992,blanch2015}. A user tracks $n$ satellites
of a lunar constellation. About a nominal linearization point, the vector of
measurement residuals $\yv\in\R^{n}$ relates to the state correction
$\xv=(\delta\bm r,\,c\,\delta t)\in\R^{4}$---three position components and one receiver
clock component---through
\begin{equation}
  \yv \;=\; \Gmat\,\xv \;+\; \bm\varepsilon,
  \qquad
  \Gmat \;=\;
  \begin{bmatrix}
    -\los_1^\top & 1 \\
    \vdots & \vdots \\
    -\los_n^\top & 1
  \end{bmatrix}
  \in \R^{n\times4},
  \label{eq:measmodel}
\end{equation}
where $\los_i=(\sv_i-\uv)/\|\sv_i-\uv\|$ is the unit line-of-sight (LOS) vector from the
user at $\uv$ to satellite $i$ at $\sv_i$, and $\bm\varepsilon$ is measurement noise. The
$i$-th row of the geometry (design) matrix $\Gmat$ is $[-\los_i^\top,\,1]$: the position
partial is the negative LOS and the clock partial is unity. We assume $n>4$ and
$\rank(\Gmat)=4$, so the problem is redundant---the precondition for any residual-based
integrity test.

\subsection{Range space, parity space, and the two projectors}
\label{sec:projectors}

The column space $\rng(\Gmat)\subseteq\R^{n}$ is the four-dimensional subspace of
measurement vectors that a state can produce exactly; its orthogonal complement
$\rng(\Gmat)^{\perp}$, of dimension $n-4$, is the \emph{parity space}. Let
\begin{equation}
  \Ppar \;=\; \Gmat(\Gmat^\top\Gmat)^{-1}\Gmat^\top,
  \qquad
  \Pperp \;=\; \bm I_n - \Ppar,
  \label{eq:projectors}
\end{equation}
be the orthogonal projectors onto $\rng(\Gmat)$ (rank $4$) and onto the parity space
(rank $n-4$), respectively. A parity matrix $\Smat\in\R^{(n-4)\times n}$ with
$\Smat\Gmat=\bm 0$ and $\Smat\Smat^\top=\bm I_{n-4}$ realizes the same complement, so
that $\Pperp=\Smat^\top\Smat$ and the \emph{parity vector} is $\bm p=\Smat\yv$. The
least-squares state estimate and its residual are
\begin{equation}
  \hat\xv \;=\; (\Gmat^\top\Gmat)^{-1}\Gmat^\top\yv,
  \qquad
  \bm r \;=\; \yv - \Gmat\hat\xv \;=\; \Pperp\yv,
  \qquad
  \|\bm r\| \;=\; \|\Smat\yv\| \;=\; \|\bm p\|.
  \label{eq:residual}
\end{equation}
The snapshot RAIM/ARAIM fault-detection statistic is the parity-vector norm
$\|\bm p\|=\|\Pperp\yv\|$ (equivalently the sum of squared residuals): a fault is
declared when it exceeds a threshold set by the noise model and the allowed false-alarm
rate \citep{sturza1988,brown1992,parkinson1988}.

\subsection{Detectability lives only in the parity projection}
\label{sec:detectability}

Consider an additive measurement error $\delta\yv$ superimposed on the nominal
measurement. Only its parity projection changes the detection statistic, and only its
range projection changes the state estimate:
\begin{equation}
  \bm p(\yv+\delta\yv) - \bm p(\yv) \;=\; \Smat\,\delta\yv \;=\; \Smat\,\Pperp\,\delta\yv,
  \qquad
  \hat\xv(\yv+\delta\yv) - \hat\xv(\yv) \;=\; (\Gmat^\top\Gmat)^{-1}\Gmat^\top\delta\yv .
  \label{eq:split}
\end{equation}
We record the three engine quantities that Sections~\ref{sec:anchor}--\ref{sec:cmpl}
use, all functions of a single injected error $\delta\yv$:
\begin{align}
  \text{blind state}\quad
    &\code{blind\_dx} = (\Gmat^\top\Gmat)^{-1}\Gmat^\top\delta\yv \in\R^4,
    \label{eq:blinddx}\\
  \text{blind fraction}\quad
    &\code{blind\_fraction} = \frac{\|\Ppar\,\delta\yv\|}{\|\delta\yv\|} \in[0,1],
    \label{eq:blindfrac}\\
  \text{detectable norm}\quad
    &\code{detectable\_norm} = \|\Pperp\,\delta\yv\| = \|\Smat\,\delta\yv\| .
    \label{eq:detnorm}
\end{align}
The blind fraction is the share of $\delta\yv$ that the residual test cannot see (the
fraction absorbed into the state), and the detectable norm is the induced change in the
detection statistic. The following elementary but load-bearing statement is classical
and we claim no novelty for it; it is the reference point the rest of the paper builds
on.

\begin{proposition}[Parity-space blindness; classical]
\label{prop:blind}
A measurement error $\delta\yv$ is undetectable by the snapshot residual test if and
only if it lies in $\rng(\Gmat)$, in which case $\Pperp\delta\yv=\bm 0$
(\code{detectable\_norm}$=0$), $\code{blind\_fraction}=1$, and $\delta\yv$ is absorbed
without residual into the state error $\hat\xv$ of \eqref{eq:blinddx}. Conversely any
component of $\delta\yv$ in $\rng(\Gmat)^{\perp}$ enters the statistic in full.
\end{proposition}

\begin{proof}
$\delta\yv\in\rng(\Gmat)$ means $\delta\yv=\Gmat\bm a$ for some $\bm a\in\R^4$; then
$\Pperp\delta\yv=(\bm I-\Ppar)\Gmat\bm a=\Gmat\bm a-\Gmat\bm a=\bm 0$ by
$\Ppar\Gmat=\Gmat$, so the statistic is unchanged and, by \eqref{eq:blinddx},
$\hat\xv$ shifts by $(\Gmat^\top\Gmat)^{-1}\Gmat^\top\Gmat\bm a=\bm a$. If instead
$\Pperp\delta\yv\ne\bm 0$ then $\|\bm p\|$ changes by $\|\Smat\delta\yv\|>0$. The two
subspaces are orthogonal and complementary, so the decomposition and the equivalence
are exact.
\end{proof}

\noindent Proposition~\ref{prop:blind} is stated in \citet{sturza1988},
\citet{brown1992}, and \citet{parkinson1988}, and it is the property our own detector
records in its \code{parity\_raim\_test}. Everything that is genuinely ours starts one
step later: identifying \emph{which physically meaningful constellation errors} land in
this blind subspace (Section~\ref{sec:theorem}), measuring their real magnitude
(Section~\ref{sec:anchor}), and bounding the specifically multi-provider part that no
single provider's ARAIM budget covers (Section~\ref{sec:cmpl}).

\section{The rigid-body constellation error is blind}
\label{sec:theorem}

This section states and proves the paper's proven contribution: that the full
rigid-body transformation of the constellation lies in $\rng(\Gmat)$, and hence
(Proposition~\ref{prop:blind}) is invisible to the snapshot residual test. We then make
the connection---a synthesis, not a discovery---to the geodetic datum defect of the
companion paper \citep{baweja2026datum}. The algebra is \Proven{}; the physical
interpretation and every later magnitude are labelled where they appear.

\subsection{The rigid-body lemma}
\label{sec:rigidlemma}

Let the constellation undergo a common rigid-body transformation: a translation
$\bm\tau\in\R^3$ and an infinitesimal rotation $\bm\theta\in\R^3$ about a pivot $O$,
applied identically to every satellite,
\begin{equation}
  \ds_i \;=\; \bm\tau \;+\; \bm\theta\times(\sv_i-O),
  \qquad i=1,\dots,n .
  \label{eq:rigid}
\end{equation}
Such a common shift is exactly what a change of the broadcast reference frame, or a
frame/ephemeris offset shared by every satellite of one provider, imposes on the
satellite positions. To first order a satellite-position error $\ds_i$ enters the
measurement residual through its LOS projection,
\begin{equation}
  \delta y_i \;=\; \los_i\cdot\ds_i,
  \label{eq:losproj}
\end{equation}
because the modelled range to satellite $i$ depends on the satellite position only
through $\|\sv_i-\uv\|$, whose gradient is $\los_i$.

\begin{lemma}[The rigid-body constellation error is blind]
\label{lem:rigid}
For any user position $\uv$ and pivot $O$, the first-order measurement error
\eqref{eq:losproj} induced by the rigid-body transformation \eqref{eq:rigid} lies in
$\rng(\Gmat)$:
\begin{equation}
  \delta\yv \;=\; \Gmat\begin{bmatrix}-(\bm\tau+\bm\theta\times(\uv-O))\\[2pt] 0\end{bmatrix}
  \;\in\;\rng(\Gmat).
  \label{eq:rigidinrange}
\end{equation}
It therefore produces a zero parity vector and is undetectable by the snapshot residual
test; it is absorbed, without residual, into a user position error
$-\bigl(\bm\tau+\bm\theta\times(\uv-O)\bigr)$ and zero clock error. Both the translation
($3$ degrees of freedom) and the rotation ($3$ degrees of freedom) are fully blind. The
smallest detectable residual of a \emph{genuine} rigid rotation is second order in the
angle, of magnitude $\mathcal{O}(|\bm\theta|^2\,L)$ for a lever arm $L$; only non-rigid
deformation of the constellation is detectable at first order.
\end{lemma}

\begin{proof}
Substitute \eqref{eq:rigid} into \eqref{eq:losproj} and split the rotation lever arm
$\sv_i-O=(\sv_i-\uv)+(\uv-O)$:
\begin{equation}
  \delta y_i
  \;=\; \underbrace{\los_i\cdot\bm\tau}_{\text{translation}}
     \;+\; \underbrace{\los_i\cdot\bigl(\bm\theta\times(\sv_i-\uv)\bigr)}_{\text{(A)}}
     \;+\; \underbrace{\los_i\cdot\bigl(\bm\theta\times(\uv-O)\bigr)}_{\text{(B)}} .
  \label{eq:threeterms}
\end{equation}
\emph{Term (A) vanishes.} By the cyclic invariance of the scalar triple product,
$\los_i\cdot(\bm\theta\times(\sv_i-\uv))=\bm\theta\cdot((\sv_i-\uv)\times\los_i)$, and
$\los_i$ is parallel to $\sv_i-\uv$ (it is the corresponding unit vector), so the cross
product $(\sv_i-\uv)\times\los_i=\bm 0$. This is the crux: the rotation moves each
satellite \emph{perpendicular} to nothing the receiver can sense along its own
sightline; the LOS projection of the along-arc motion is a scalar triple product with a
repeated vector and is identically zero.

\emph{Term (B) is a common constant.} Set $\bm a:=\bm\theta\times(\uv-O)$, which does
\emph{not} depend on the satellite index $i$. Then term (B) is $\los_i\cdot\bm a$, and
the translation term is $\los_i\cdot\bm\tau$; combining, $\delta y_i=\los_i\cdot(\bm\tau
+\bm a)=-\los_i\cdot\bigl(-(\bm\tau+\bm a)\bigr)$. But the $i$-th row of $\Gmat$ acting on
a state $(\bm b,0)$ gives exactly $-\los_i\cdot\bm b$, so with $\bm b=-(\bm\tau+\bm a)$
we obtain \eqref{eq:rigidinrange}: $\delta\yv=\Gmat[-(\bm\tau+\bm\theta\times(\uv-O));0]$.
By Proposition~\ref{prop:blind}, $\Pperp\delta\yv=\bm 0$ and the state absorbs
$-(\bm\tau+\bm a)$.

\emph{User-independence and the second-order remark.} The membership
$\delta\yv\in\rng(\Gmat)$ holds for every $\uv$ and every pivot $O$; only the value of
the absorbed state $\bm a=\bm\theta\times(\uv-O)$ depends on where the user is. That the
error is \emph{undetectable} is thus a property of the transformation, not of the user
geometry. Finally, \eqref{eq:losproj} is the first-order (linearized) map; for a genuine
finite rotation $\bm R(\bm\theta)$ the exact along-arc displacement of a satellite is
$\bm R(\bm\theta)(\sv_i-O)-(\sv_i-O)$, whose component not captured by the first-order
term (A) is $\mathcal{O}(|\bm\theta|^2)$. Projected on the LOS and scaled by the lever
arm $L\sim\|\sv_i-\uv\|$, the leftover parity signal is $\mathcal{O}(|\bm\theta|^2 L)$;
for nanoradian-scale frame errors ($|\bm\theta|\sim\SI{5}{\nano\radian}$) and a lever
arm $L\sim\SI{3.84e8}{m}$ this is on the order of $10^{-9}\,$m---negligible. Only genuine
\emph{non-rigid} deformation of the constellation (relative motion between satellites)
carries a first-order parity signal and is detectable.
\end{proof}

\begin{remark}[The clock degree of freedom is likewise blind]
\label{rem:clock}
The same argument shows a common receiver-clock offset absorbs the all-ones column: a
constant $\delta y_i=c\,\delta t$ for all $i$ is $\Gmat[\bm 0;c\,\delta t]\in\rng(\Gmat)$.
The four state directions---three position, one clock---are the blind subspace by
construction; Lemma~\ref{lem:rigid} adds that three \emph{rotational} constellation
degrees of freedom map into these same four state directions and are blind for the same
reason.
\end{remark}

\subsection{The datum defect is a subset of the blind subspace}
\label{sec:datumtie}

Lemma~\ref{lem:rigid} has a clean geodetic reading. The rigid-body group---three
translations and three rotations of a point set---is precisely the \emph{datum defect}
(gauge freedom) of a free-network adjustment: the family of transformations that a
network's own internal geometry cannot fix, and that must be pinned by an external
datum \citep{baarda1973,teunissen2000}. Baarda's $S$-transformation parameterizes
exactly this freedom, and the companion paper \citep{baweja2026datum} classifies its
lunar realization---which components of the single-provider datum are unobservable, and
what observation geometry (libration, an absolute tie) resolves them.

\begin{corollary}[Datum defect $\subset$ RAIM blind subspace]
\label{cor:datum}
Let $\mathcal{D}$ be the rigid-body (translation$+$rotation) datum-defect subspace of a
provider's constellation---the gauge directions of \citep{baweja2026datum} that no
internal observation fixes. Then the first-order measurement image of every
$\bm d\in\mathcal{D}$ lies in $\rng(\Gmat)$, so
$\mathcal{D}\subseteq$ the RAIM blind subspace. An error a provider's own network
cannot fix is, a fortiori, an error that provider's snapshot integrity monitor cannot
see.
\end{corollary}

\begin{proof}
Immediate from Lemma~\ref{lem:rigid}: every generator of the rigid-body group
(translation $\bm\tau$, rotation $\bm\theta$) maps under \eqref{eq:losproj} into
$\rng(\Gmat)$, and the image of a subspace under a linear map is the span of the images
of its generators.
\end{proof}

\begin{remark}[This is a synthesis, not a discovery]
\label{rem:synthesis}
We claim novelty for none of Section~\ref{sec:theorem}'s statements as isolated facts.
That a common translation of all satellites is indistinguishable from a user shift is
folklore in the GNSS integrity literature; that the datum defect is the rigid-body
gauge is standard geodesy \citep{baarda1973,teunissen2000}; and that RAIM is blind to
$\rng(\Gmat)$ is Proposition~\ref{prop:blind}. The contribution is the explicit
\emph{connection}---that the full rotational datum freedom, not only translation, is in
the blind subspace, and that the lunar datum-defect classification of P1
\citep{baweja2026datum} therefore lands entirely inside the P3 integrity blind subspace.
The subset relation of Corollary~\ref{cor:datum} is what carries into the multi-provider
budget of Section~\ref{sec:cmpl}: the very frame freedom that a single provider cannot
observe becomes, \emph{differenced across independent providers}, a real position error
that no provider's residual test detects.
\end{remark}

\section{The real inter-ephemeris common-mode}
\label{sec:anchor}

We now attach a real magnitude to the blind subspace. The section has two parts, and
their honesty status is different: Part~A is an \emph{analytic identity} exercised
numerically, and Part~B is the paper's one externally \Validated{} row. We are explicit
about what each earns.

\subsection{Part A: the analytic-identity check}
\label{sec:partA}

The common-shift construction of Lemma~\ref{lem:rigid} guarantees, by algebra, that an
error built as $\delta\yv=\Gmat[-\ds;0]$ lies in $\rng(\Gmat)$ and is fully blind.
Part~A confirms the engine reproduces this identity to machine precision. We inject a
state-equivalent common shift with position/clock signature
$\bm\delta=(1.000,\,-2.000,\,0.500,\,3.000)$ (metres, and metres of clock range), form
$\delta\yv=\Gmat\bm\delta$, and recover the state:
\begin{equation}
\begin{aligned}
  &\code{blind\_fraction}=1.000000000000000, \\
  &\max_j\bigl|\code{blind\_dx}_j-\bm\delta_j\bigr|=2.5\times10^{-15},\qquad
  \code{detectable\_norm}\approx2.7\times10^{-15}\,\text{m}.
\end{aligned}
  \label{eq:partA}
\end{equation}
The recovered blind state matches the injection to $2.5\times10^{-15}$ and the
detectable norm is $2.7\times10^{-15}\,$m. We stress the interpretation demanded by
honesty: the detectable norm here is \emph{floating-point zero}, not a small measured
residual. It is the numerical shadow of an exact algebraic identity---$\delta\yv$ was
constructed to lie in $\rng(\Gmat)$, so $\Pperp\delta\yv=\bm 0$ by
Proposition~\ref{prop:blind}, and the $\sim\!10^{-15}\,$m is the accumulated
double-precision rounding of the projection, nothing more. Part~A is an internal
consistency check, not an empirical result, and we do not present the
$\sim\!10^{-15}\,$m as a ``measured'' blindness.

\subsection{Part B: the real DE440/INPOP21a/EPM2021 magnitude}
\label{sec:partB}

The empirical content is the \emph{magnitude} of the common shift, taken from real
data. We use the geocentric-Moon position disagreement between the three independent
authoritative ephemerides---DE440 \citep{park2021de440}, INPOP21a
\citep{fienga2021inpop21a}, and EPM2021 \citep{pitjeva2021epm2021}---the same real
quantity decomposed in the companion interoperability paper
\citep{baweja2026interop}. For each provider pair and each of $366$ epochs over
2024--2025 (two-day cadence, ICRF/J2000, TDB; Appendix~\ref{app:methods}) we take the
per-epoch geocentric-Moon difference as the common frame/ephemeris shift $\ds$ seen by
the whole constellation, form $\delta\yv=\Gmat[-\ds;0]$ through the modelled lunar
geometry, and compute the parity/least-squares split \eqref{eq:blinddx}--\eqref{eq:detnorm}.
Table~\ref{tab:anchor} reports the medians over all epochs; the record count is $732$
($2$ provider pairs $\times366$ epochs).

\begin{table}[t]
\centering
\caption{The real inter-ephemeris common-mode absorbed as user position error
(medians over $366$ epochs, 2024--2025 geocentric Moon). The blind position magnitude
$|\code{blind\_dx}_{\mathrm{pos}}|$ is the empirical, \Validated{} content---the
$\approx\!2\,$m user-position error the residual test does not see. The blind fraction
($=1$) and the detectable norm ($\sim\!10^{-15}\,$m) are the \emph{analytic identity}
of the common-shift construction (Section~\ref{sec:partA}), not measurements: the
detectable norm is floating-point zero. The engine's split reproduces an independent
NumPy computation to relative and absolute error below $10^{-3}$
(Section~\ref{sec:validation}).}
\label{tab:anchor}
\small
\begin{tabular}{@{}lccc@{}}
\toprule
\textbf{Provider pair} & $|\code{blind\_dx}_{\mathrm{pos}}|$ (m) &
\code{blind\_fraction} & \code{detectable\_norm} (m) \\
 & \textit{(empirical magnitude)} & \textit{(analytic identity)} & \textit{(f.p.\ zero)} \\
\midrule
DE440--INPOP21a & $2.3955$ & $1.000000000$ & $\sim2\times10^{-15}$ \\
DE440--EPM2021  & $2.0050$ & $1.000000000$ & $\sim2\times10^{-15}$ \\
\bottomrule
\end{tabular}
\end{table}

The median blind position magnitude is $\SI{2.3955}{m}$ for DE440--INPOP21a and
$\SI{2.0050}{m}$ for DE440--EPM2021; the median blind fraction is
$1.000000000$ and the median detectable norm $\sim\!2\times10^{-15}\,$m for every one of
the $366$ epochs of both pairs. The reading is: a real $\approx\!2\,$m inter-ephemeris
frame/ephemeris disagreement, applied as a constellation common-mode, is absorbed
\emph{almost entirely} into user position error and is structurally invisible to the
RAIM/ARAIM residual test.

\subsection{The honesty boundary: what realness contributes}
\label{sec:honestyboundary}

It would be easy---and wrong---to read Table~\ref{tab:anchor} as ``the data show the
error is $100\%$ blind.'' The blindness is \emph{not} a finding of the data. The blind
fraction is exactly $1$ and the detectable norm exactly zero (to machine precision)
because we \emph{constructed} $\delta\yv=\Gmat[-\ds;0]$ to lie in $\rng(\Gmat)$; by
Proposition~\ref{prop:blind} this is forced, whatever the numerical value of $\ds$. What
the real ephemerides contribute is one thing only: the \emph{magnitude}
$\|\ds\|\approx2\,$m, i.e.\ how large the blind user-position error actually is for a
plausible inter-provider frame/ephemeris disagreement. Equivalently: the \emph{fact} of
blindness is proven (Section~\ref{sec:theorem}); the \emph{size} of the blind error is
what the DE440/INPOP21a/EPM2021 differences pin down. The externally \Validated{} claim
is correspondingly narrow and is the subject of Section~\ref{sec:validation}: the
engine's Rust parity/least-squares split reproduces an independent NumPy computation, on
these real position magnitudes, to relative and absolute error below $10^{-3}$. We
validate the linear algebra at real magnitudes; we do not, and cannot, ``validate'' an
identity.

\section{The cross-provider protection level and the integrity envelope}
\label{sec:cmpl}

We now turn the blind magnitude of Section~\ref{sec:anchor} into an integrity bound.
Every quantity in this section is \Modelled{} and tagged \modtag{}: the constellation
geometry, the common-mode covariance, the cross-provider protection level (CMPL), the
representative ARAIM horizontal protection level (HPL), and the total envelope are all
representative---the ``providers'' are ephemerides mapped through a modelled
constellation, and \emph{no two lunar-PNT providers fly}. The section's purpose is to
state the load-bearing claim in its correct form and to exhibit the envelope's
structure, not to certify a number.

\subsection{Why a new term is needed: the load-bearing claim, stated correctly}
\label{sec:loadbearing}

The tempting claim---``ARAIM is provably blind to common-mode''---is false, and a
reviewer is right to reject it. The distinction is between the residual test and the
protection level.

\begin{itemize}
  \item The RAIM/ARAIM \emph{residual test} is provably blind to $\rng(\Gmat)$
  (Proposition~\ref{prop:blind}, Lemma~\ref{lem:rigid}). A common-mode frame/ephemeris
  error triggers no fault flag.

  \item The ARAIM \emph{protection level} does \emph{not} rely on the residual test to
  catch this error. The multiple-hypothesis solution-separation (MHSS) protection level
  bounds the position error under the \emph{fault-free} (nominal, $H_0$) hypothesis
  through the user range accuracy overbound $\sigma_{\mathrm{URA}}$ and the nominal bias
  $b_{\mathrm{nom}}$ declared in the integrity support message (ISM)
  \citep{blanch2015,wgc2016araim,decleene2000}. A residual-invisible error is therefore
  still budgeted---up to whatever the ISM declares. Within a single provider, an
  ephemeris/frame bias that the residual test cannot see is covered by
  $\sigma_{\mathrm{URA}}$ and $b_{\mathrm{nom}}$, provided the ISM overbound is honest
  \citep{montenbruck2018sisre}.
\end{itemize}

The genuinely uncovered term is specific to the \emph{multi-provider} case. When a user
fuses, or hands off between, two \emph{independent} lunar providers, their mutual
frame/ephemeris common-mode is not contained in \emph{either} provider's per-provider
URA. Each provider's $\sigma_{\mathrm{URA}}$ overbounds \emph{its own} signal-in-space
error against \emph{its own} realized frame and ephemeris; it says nothing about the
offset between two independent realizations. The measured inter-ephemeris
common-mode---$\approx\!2\,$m between DE440, INPOP21a, and EPM2021
(Section~\ref{sec:anchor}, and the interoperability floor of \citep{baweja2026interop})%
---is exactly such a cross-provider term. It is residual-invisible (Lemma~\ref{lem:rigid})
\emph{and} outside any single provider's ISM budget. That is the gap the CMPL fills.

\subsection{The common-mode covariance and the CMPL}
\label{sec:cmpldef}

We model the cross-provider common shift as a zero-mean Gaussian with an isotropic
$3$-D covariance whose total one-sigma magnitude is fixed to the Part-B floor,
\begin{equation}
  \Sig_{\mathrm{cm}} \;=\; \frac{\sigma_{\mathrm{cm}}^2}{3}\,\bm I_3,
  \qquad
  \sigma_{\mathrm{cm}} \;=\; \|\ds\| \;=\; \SI{2.0000}{m}\ \text{(1-}\sigma\text{)},
  \label{eq:cmcov}
\end{equation}
sized to the $\approx\!2\,$m inter-ephemeris floor measured in Section~\ref{sec:anchor}
and in the companion interoperability paper \citep{baweja2026interop}, which is the
input that feeds $\Sig_{\mathrm{cm}}$. Because a common shift maps (Lemma~\ref{lem:rigid})
directly into user position error, the blind position covariance is $\Sig_{\mathrm{cm}}$
itself. Projecting onto the horizontal plane gives a horizontal one-sigma
\begin{equation}
  \sigma_H \;=\; \sqrt{\tfrac{2}{3}}\,\sigma_{\mathrm{cm}} \;=\; \SI{1.6330}{m},
  \label{eq:sigmaH}
\end{equation}
(two of three isotropic axes are horizontal), and the cross-provider protection level is
the standard $2$-D horizontal Rayleigh inflation of this one-sigma,
\begin{equation}
  \mathrm{CMPL} \;=\; K_H\,\sigma_H \;=\; 6.0\times\SI{1.6330}{m} \;=\; \SI{9.7980}{m},
  \qquad K_H = 6.0,
  \label{eq:cmpl}
\end{equation}
with $K_H=6.0$ the representative horizontal integrity multiplier (the $2$-D Rayleigh
factor used at civil-aviation integrity risks in ARAIM \citep{blanch2015,wgc2016araim}).
We label \eqref{eq:cmpl} plainly: the CMPL value is an \emph{input-echo mechanism
illustration}, not a finding. It restates the $\approx\!2\,$m modelled input through a
fixed multiplier; its role is to show \emph{how} a cross-provider term enters an
integrity budget and with what sign and magnitude, not to certify $\SI{9.80}{m}$ as a
lunar CMPL.

\subsection{The integrity envelope, and why the sum is a triangle bound}
\label{sec:envelope}

We combine the CMPL with a representative single-provider ARAIM horizontal protection
level. Taking $\mathrm{HPL}_{\mathrm{ARAIM}}=\SI{15.0}{m}$ (representative and
illustrative---\emph{not} \Validated{}; a placeholder for whatever a real lunar ARAIM
configuration would compute), the total cross-provider integrity envelope is
\begin{equation}
  \mathrm{HPL}_{\mathrm{total}}
  \;=\; \mathrm{HPL}_{\mathrm{ARAIM}} + \mathrm{CMPL}
  \;=\; \SI{15.0}{m} + \SI{9.7980}{m}
  \;=\; \SI{24.7980}{m}.
  \label{eq:hpltotal}
\end{equation}
Table~\ref{tab:envelope} collects the terms. Two properties of \eqref{eq:hpltotal} are
essential and both are stated as \Modelled{} guidance.

\begin{table}[t]
\centering
\caption{The cross-provider integrity envelope (\Modelled{}, \modtag{}). The
common-mode one-sigma $\sigma_{\mathrm{cm}}=\SI{2}{m}$ is sized to the real
inter-ephemeris floor (Section~\ref{sec:anchor}); the CMPL is an input-echo mechanism
illustration, not a finding. The total is a \emph{conservative} triangle-inequality sum:
a root-sum-square ($\sqrt{15^2+9.798^2}=\SI{17.92}{m}$) would under-bound and overclaim
integrity, because the ARAIM and common-mode contributions can align in the position
domain. The envelope is valid only where the common-mode is not already absorbed into
$\sigma_{\mathrm{URA}}$ (Section~\ref{sec:conditionality}).}
\label{tab:envelope}
\small
\begin{tabular}{@{}lcl@{}}
\toprule
\textbf{Term} & \textbf{Value (m)} & \textbf{Status} \\
\midrule
$\sigma_{\mathrm{cm}}=\|\ds\|$ (common-mode 1-$\sigma$) & $2.0000$ & \Modelled{} input, sized to real floor \\
$\sigma_H=\sqrt{2/3}\,\sigma_{\mathrm{cm}}$ (horizontal 1-$\sigma$)   & $1.6330$ & \Modelled{} \\
$\mathrm{CMPL}=K_H\sigma_H$, $K_H=6.0$    & $9.7980$ & \Modelled{} (input-echo illustration) \\
$\mathrm{HPL}_{\mathrm{ARAIM}}$ (representative)          & $15.0000$ & \Modelled{}, \emph{not} \Validated{} \\
$\mathrm{HPL}_{\mathrm{total}}=\mathrm{HPL}_{\mathrm{ARAIM}}+\mathrm{CMPL}$ & $24.7980$ & \Modelled{} (conservative triangle sum) \\
\midrule
(root-sum-square, for contrast---\emph{under-bounds}) & $17.92$ & rejected: overclaims integrity \\
\bottomrule
\end{tabular}
\end{table}

\paragraph{The sum is a triangle bound, not a root-sum-square.}
Integrity budgets normally combine independent error sources in quadrature. Here that
would be wrong. The single-provider ARAIM error and the cross-provider common-mode are
\emph{not} guaranteed independent in the position domain: their horizontal
contributions can align, so the worst-case combined displacement is the arithmetic sum,
not the root-sum-square. Using $\sqrt{\mathrm{HPL}_{\mathrm{ARAIM}}^2+\mathrm{CMPL}^2}
=\SI{17.92}{m}$ would under-bound the true worst case of $\SI{24.80}{m}$ and
\emph{overclaim} integrity---the cardinal sin of an integrity monitor. The
triangle-inequality sum \eqref{eq:hpltotal} is the conservative choice and the correct
one for a bound.

\subsection{The anti-double-counting conditionality}
\label{sec:conditionality}

The envelope \eqref{eq:hpltotal} is valid \emph{only where the common-mode is not
already absorbed into $\sigma_{\mathrm{URA}}$}. This condition is not a footnote; it is
the boundary of the CMPL's validity.

\begin{itemize}
  \item If each provider's ISM declares an \emph{un-inflated} $\sigma_{\mathrm{URA}}$
  that covers only its own signal-in-space error against its own realization---the
  normal case---then the inter-provider common-mode is genuinely uncovered, and adding
  the CMPL is necessary. The envelope \eqref{eq:hpltotal} then restores a valid bound.

  \item If, instead, a provider inflates $\sigma_{\mathrm{URA}}$ (or $b_{\mathrm{nom}}$)
  to \emph{also} cover the cross-provider disagreement, then the common-mode is already
  inside the ARAIM fault-free term, and adding the CMPL \emph{double-counts}: the true
  bound is already $\mathrm{HPL}_{\mathrm{ARAIM}}$ alone, and
  $\mathrm{HPL}_{\mathrm{ARAIM}}+\mathrm{CMPL}$ is merely conservative-to-the-point-of-%
  waste, not wrong, but no longer tight.
\end{itemize}

Hence the precise statement:
$\mathrm{HPL}_{\mathrm{total}}=\mathrm{HPL}_{\mathrm{ARAIM}}+\mathrm{CMPL}$ is the valid
bound for the \emph{un-inflated-URA} case---the case where a provider budgets only its
own realization and the cross-provider term is left to a multi-provider integrity layer.
The design consequence for a multi-provider architecture is a clean choice: either
mandate that each provider's ISM overbound explicitly \emph{includes} the permissible
inter-provider common-mode (folding CMPL into $\sigma_{\mathrm{URA}}$ upstream), or
carry a separate CMPL term downstream at the fusion point. What must not happen is
\emph{neither}---leaving the cross-provider common-mode in no budget at all---which is
exactly the gap Lemma~\ref{lem:rigid} and Section~\ref{sec:anchor} expose: it is
residual-invisible \emph{and}, absent an explicit provision, outside every
per-provider URA.

\section{Validation and the structure/magnitude firewall}
\label{sec:validation}

A paper that converts a blind constellation error into an integrity design rule must be
scrupulous about which numbers are earned on real data and which are representative. We
enforce this with an explicit firewall, implemented as a rule in the engine's
verification suite and stated plainly here: a quantity is \Validated{} only when it is
checked against an \emph{independent} oracle---a separate library, dataset, or published
vector---to a stated tolerance; a result is \Proven{} when it is an algebraic theorem;
a fact taken from the literature is \emph{cited}; and everything else is \Modelled{}.

\subsection{The independent oracle for the split}
\label{sec:oracle}

The one load-bearing empirical object of this paper is the parity/least-squares split of
the real inter-ephemeris common-mode at real magnitudes (Table~\ref{tab:anchor}). We
validate it against an independent NumPy reference. A generator script
(Section~\ref{sec:availability}) reads the DE440/INPOP21a/EPM2021 kernels through the
IMCCE \code{calceph} library, samples the geocentric Moon on the same grid used by the
engine, forms the per-epoch common shift $\ds$, builds the same modelled constellation
geometry $\Gmat$, and computes the projectors \eqref{eq:projectors} and the three split
quantities \eqref{eq:blinddx}--\eqref{eq:detnorm} by an independent NumPy linear-algebra
path (an explicit QR/normal-equations solve rather than the engine's routine), emitting
the expected \code{blind\_dx}, \code{blind\_fraction}, and \code{detectable\_norm} to a
vendored fixture (\code{reference.json}). The engine's split reproduces every value to a
relative error below $10^{-3}$ and an absolute error below $\SI{e-3}{m}$ on the
metre-scale position magnitudes. This is the paper's single \Validated{} row: the
\emph{linear algebra} of the split, on the \emph{real} ephemeris magnitudes, is
externally confirmed.

We restate the boundary made precise in Section~\ref{sec:honestyboundary}. What is
validated is that two independent implementations agree on the metre-scale blind
position error. What is \emph{not} an empirical claim---and is not presented as one---is
the blindness itself: $\code{blind\_fraction}=1$ and $\code{detectable\_norm}\approx0$
are analytic identities of the construction $\delta\yv=\Gmat[-\ds;0]$
(Proposition~\ref{prop:blind}), true to machine precision for any $\ds$ whatsoever. The
$\sim\!10^{-15}\,$m detectable norm is floating-point zero and is never quoted as a
measurement.

\subsection{What is cited, proven, validated, and modelled}
\label{sec:matrix}

Table~\ref{tab:honesty} draws the line explicitly.

\begin{table}[t]
\centering
\caption{The structure/magnitude firewall. Only the parity/least-squares split at real
magnitudes is \Validated{} (independent NumPy oracle on real DE440/INPOP21a/EPM2021
data). The parity-space blindness is classical prior art (\emph{cited}); the rigid-body
lemma and the datum-defect$\,\subset\,$blind relation are \Proven{}; the analytic
identity is an internal check; every system-level magnitude is \Modelled{}.}
\label{tab:honesty}
\small
\begin{tabular}{@{}p{5.2cm}p{3.6cm}p{2.6cm}c@{}}
\toprule
\textbf{Result} & \textbf{Basis} & \textbf{Oracle} & \textbf{Status} \\
\midrule
Parity-space blindness of $\rng(\Gmat)$ (Prop.~\ref{prop:blind}) &
Sturza 1988; Brown 1992; Parkinson \& Axelrad 1988 &
--- (prior art) & \emph{cited} \\[3pt]
Rigid-body error $\in\rng(\Gmat)$; datum defect $\subset$ blind subspace
(Lem.~\ref{lem:rigid}, Cor.~\ref{cor:datum}) &
algebraic proof &
--- & \Proven{} \\[3pt]
Split of the real inter-ephemeris common-mode at $\approx\!2\,$m
(Table~\ref{tab:anchor}, Part~B) &
real DE440/INPOP/EPM vs NumPy, $<\!10^{-3}$ &
ext.\ library + real data & \Validated{} \\[3pt]
Analytic-identity check (Part~A, Eq.~\ref{eq:partA}) &
common-shift construction $\delta\yv=\Gmat[-\ds;0]$ &
internal (f.p.\ zero) & identity \\[3pt]
Common-mode covariance $\Sig_{\mathrm{cm}}$; CMPL; total envelope
(Table~\ref{tab:envelope}) &
Rayleigh inflation of modelled 1-$\sigma$ &
internal & \Modelled{} \\[3pt]
Representative $\mathrm{HPL}_{\mathrm{ARAIM}}=\SI{15}{m}$ &
illustrative placeholder &
--- & \Modelled{} \\
\bottomrule
\end{tabular}
\end{table}

Two internal checks are honest but not independent, and are labelled so: Part~A is an
analytic identity (a synthetic common shift is injected and its blind fraction confirmed
at $1$), and the envelope assembly is an internal combination of the modelled
common-mode one-sigma with a representative HPL. Neither is claimed as externally
validated.

\subsection{Reproducibility}
\label{sec:reproducibility}

Every quantitative result---Part~A's identity check, Part~B's real split
(Table~\ref{tab:anchor}), and the envelope of Table~\ref{tab:envelope})---is
regenerated from the committed engine and the vendored fixtures by a single command
(Section~\ref{sec:availability}). The \Validated{} row is asserted in the
continuous-integration test suite, which enforces the external-oracle gate: the engine's
split must match \code{reference.json} to the stated tolerance or the build fails. There
are no hand-entered numbers in the tables; the manuscript values and the engine outputs
are the same numbers, traceable to one run.

\section{Discussion}
\label{sec:discussion}

\subsection{Positioning against the RAIM/ARAIM and ISM literature}
\label{sec:positioning}

We place the result precisely against prior art, in five statements that also fix, once
more, exactly what is and is not new.

First, \emph{parity-space blindness is classical} \citep{sturza1988,brown1992,%
parkinson1988} and is even noted in our own detector---no novelty is claimed for it. The
snapshot residual test sees a fault only through the parity subspace, and an error in
$\rng(\Gmat)$ has no parity signature; this is the starting point, not the contribution.

Second, \emph{a rigid frame transform common to the satellites (and user-independent) is
the geodetic datum-defect / gauge freedom, wholly in $\rng(\Gmat)$}
\citep{baarda1973,teunissen2000}; we make explicit that the P1 single-provider
datum-defect subspace \citep{baweja2026datum} is a subset of the RAIM blind subspace
(Corollary~\ref{cor:datum}). This is a synthesis of two known facts---the rigid-body
gauge of free-network geodesy and the range-space blindness of RAIM---not a new
phenomenon; its value is that it identifies \emph{which} physically meaningful error
(the full three-rotation frame freedom, not just translation) is blind, and ties the
integrity blindness to the identifiability classification of the companion paper.

Third, \emph{ARAIM does not ignore undetectable errors}: the MHSS protection level
\citep{blanch2015,wgc2016araim} bounds the fault-free ($H_0$) hypothesis via the
$\sigma_{\mathrm{URA}}$ overbound and the nominal bias $b_{\mathrm{nom}}$, so
residual-invisible errors are still budgeted up to what the ISM declares
\citep{decleene2000,montenbruck2018sisre}. This is why the naive claim---``ARAIM is
blind to common-mode''---is wrong at the protection-level layer, and we are careful to
reject it.

Fourth, \emph{the uncovered term is specific to the lunar multi-provider case}: an
inter-provider frame/ephemeris common-mode of $\approx\!2\,$m (measured on the real
DE440/INPOP21a/EPM2021 differences, Section~\ref{sec:anchor}) is not contained in any
single provider's per-provider URA, because each URA overbounds one provider's error
against its own realization, not the offset between two independent realizations. This is
the genuine gap, and it exists only once more than one independent provider is fused.

Fifth, \emph{the CMPL adds to the ARAIM HPL only to the extent the common-mode exceeds
the $\sigma_{\mathrm{URA}}$ overbound}: where $\sigma_{\mathrm{URA}}$ is inflated to cover
it, the CMPL is already inside the ARAIM fault-free term and adding it double-counts, so
$\mathrm{HPL}_{\mathrm{total}}=\mathrm{HPL}_{\mathrm{ARAIM}}+\mathrm{CMPL}$ is a
conservative bound for the un-inflated-URA case (Section~\ref{sec:conditionality}). The
CMPL is not a blanket surcharge; it is the term that closes a specifically
cross-provider, specifically un-budgeted gap.

\paragraph{Honest-novelty statement, restated.}
The blindness of snapshot RAIM/ARAIM to errors in $\rng(\Gmat)$ is a classical property
of the parity-space formulation and is not claimed as new. Our contribution is (i) the
real-data quantification that the actual median $2.0$--$2.4\,$m inter-ephemeris
disagreement between independent authoritative lunar ephemerides (DE440 vs.\ INPOP21a
and EPM2021), projected through a lunar constellation, is absorbed almost entirely as
user position error (blind fraction $\approx1$, parity residual $\sim\!10^{-15}\,$m) and
is therefore structurally invisible to the RAIM/ARAIM residual test; (ii) the
observation that in a multi-provider lunar PNT architecture this inter-provider
common-mode lies outside any single provider's per-provider URA/ISM budget; and (iii)
the CMPL, an additive position-domain term that restores a valid bound for exactly that
uncovered cross-provider case. We do not claim novelty for the subspace blindness, for
the datum-defect$\,\subset\,$blind relation, or for bounding an undetectable error by a
protection level---the last being the mechanism of ARAIM's own fault-free/nominal-bias
term.

\subsection{What a LunaNet or Moonlight specification should require}
\label{sec:spec}

The design consequence is one requirement, stated as \Modelled{} guidance rather than a
certified budget. A multi-provider lunar PNT architecture---LunaNet and its
interoperability specification, ESA Moonlight/LCNS, or their national counterparts
\citep{israel2020lunanet,fienga2024moonlight}---must place the inter-provider
frame/ephemeris common-mode into \emph{some} explicit budget, because
Lemma~\ref{lem:rigid} shows the residual test never will and Section~\ref{sec:conditionality}
shows a per-provider URA need not either. Two admissible mechanisms follow, and the
specification should mandate one of them: \emph{(a) upstream}, require each provider's
ISM overbound $\sigma_{\mathrm{URA}}$ to explicitly include the permissible
inter-provider common-mode (so the CMPL is folded into the declared budget and the
single-provider HPL already bounds it); or \emph{(b) downstream}, define a
cross-provider integrity layer that carries a CMPL term at the fusion/hand-off point,
combined with the per-provider HPL by the conservative triangle sum \eqref{eq:hpltotal}.
What the specification must forbid is the default of \emph{neither}, which leaves a
residual-invisible, un-budgeted $\approx\!2\,$m position error in a safety-of-life
architecture. This dovetails with the companion interoperability design law
\citep{baweja2026interop}---standardize the ephemeris, not just the format---from the
integrity side: bounding the ephemeris divergence is what makes the CMPL small, and a
common designated ephemeris is what drives it to zero.

\subsection{Limitations}
\label{sec:limitations}

The firewall of Section~\ref{sec:validation} is also a statement of limitations; we make
them explicit.

\begin{itemize}
  \item \textbf{The providers are ephemerides.} The single most important caveat: no two
  lunar-PNT providers have flown, so the three ephemerides are a \emph{representative
  analog} for a real multi-provider constellation, and the constellation geometry through
  which we project is modelled. Only the parity/least-squares split at real magnitudes
  (Table~\ref{tab:anchor}) is \Validated{}; the covariance, the CMPL, the representative
  HPL, and the envelope (Table~\ref{tab:envelope}) are \Modelled{}. A real provider
  constellation would disagree for additional reasons---broadcast-ephemeris fit error,
  clock and signal-in-space errors, regional frame ties---absent from the ephemeris
  comparison.

  \item \textbf{Blindness is an identity, not a measurement.} The blind fraction of $1$
  and the $\sim\!10^{-15}\,$m detectable norm are analytic consequences of the
  common-shift construction (Proposition~\ref{prop:blind}), not empirical findings. The
  empirical content is only the $\approx\!2\,$m magnitude. We have stated this at every
  occurrence and repeat it here.

  \item \textbf{The CMPL is an input-echo illustration.} The CMPL value $\SI{9.80}{m}$ is
  a fixed Rayleigh multiplier applied to a modelled $2\,$m one-sigma; it demonstrates the
  mechanism and sign of the cross-provider term, not a certified lunar CMPL. The
  representative $\mathrm{HPL}_{\mathrm{ARAIM}}=\SI{15}{m}$ is likewise a placeholder, not
  a computed lunar protection level.

  \item \textbf{Snapshot, first-order, single-epoch.} The analysis is the standard
  snapshot linearization; it does not address time-filtered or carrier-smoothed
  monitors, multi-epoch fault growth, or the second-order non-rigid residual beyond the
  $\mathcal{O}(|\bm\theta|^2)$ estimate of Lemma~\ref{lem:rigid}. A common-mode that
  \emph{drifts} would eventually leave a detectable signature in a filtered monitor; the
  snapshot blindness is exact only for a static common shift within the epoch.

  \item \textbf{No flight data, no TRL claim.} Nothing here is a certified integrity
  monitor, a reproduction of any agency's ARAIM configuration, or a statement at
  technology readiness level above analysis. The deliverable is a cited classical
  property, a proven lemma, one validated real-magnitude split, and a modelled integrity
  envelope built on them.
\end{itemize}

None of these caveats touches the proven content---the rigid-body lemma and the
datum-defect$\,\subset\,$blind relation---or the validated content---the real-magnitude
split. Each marks where a real multi-provider deployment must supply flight data before
any absolute CMPL number is trusted.

\section{Conclusion}
\label{sec:conclusion}

Snapshot RAIM and ARAIM will be the real-time integrity monitors of multi-provider
lunar PNT, and there is one error class their residual test structurally cannot see.
That the test is blind to $\rng(\Gmat)$ is classical \citep{sturza1988,brown1992,%
parkinson1988}, and we claim no novelty for it. What we prove is that the \emph{entire}
rigid-body transformation of the constellation---three translations and three rotations,
independent of the user's own position---lies in $\rng(\Gmat)$ and is therefore fully
blind, with only second-order ($\sim\!10^{-9}\,$m) residual for realistic nanoradian
frame errors; and we make explicit, as a synthesis of free-network geodesy and RAIM,
that the single-provider datum-defect subspace of the companion paper
\citep{baweja2026datum} is a subset of this blind subspace. On the only real ``multiple
providers'' that exist today---the independent authoritative ephemerides DE440,
INPOP21a, and EPM2021---we \Validated{}, against an independent NumPy oracle to a
tolerance below $10^{-3}$, that the engine's parity/least-squares split reproduces the
absorbed $\approx\!2\,$m user-position error at real magnitudes; and we are scrupulous
that the blindness itself is an analytic identity of the common-shift construction, the
data setting only the metre magnitude, so the $\sim\!10^{-15}\,$m parity residual is
floating-point zero and never a measurement.

The load-bearing claim is stated in its correct form. ARAIM is \emph{not} provably blind
to common-mode at the protection level: its MHSS fault-free hypothesis already bounds
residual-invisible error through $\sigma_{\mathrm{URA}}$ and the nominal bias
$b_{\mathrm{nom}}$ \citep{blanch2015,wgc2016araim}. The genuinely uncovered term is the
multi-provider one---a frame/ephemeris common-mode between independent providers that
lies outside any single provider's per-provider URA/ISM budget---and the cross-provider
protection level (CMPL) restores a valid bound for exactly that case. The total envelope
$\mathrm{HPL}_{\mathrm{total}}=\mathrm{HPL}_{\mathrm{ARAIM}}+\mathrm{CMPL}$ is a
conservative triangle-inequality sum (a root-sum-square would under-bound and overclaim
integrity), valid only where the common-mode is not already absorbed into
$\sigma_{\mathrm{URA}}$, else it double-counts. The constellation, the covariance, the
CMPL, and the envelope are \Modelled{}, because the providers are ephemerides and no two
lunar-PNT providers fly.

The design rule is compact: \emph{a multi-provider lunar integrity architecture must
place the inter-provider frame/ephemeris common-mode into an explicit budget---upstream
in each provider's URA, or downstream in a cross-provider protection level---because the
residual test never will and a per-provider URA need not.} This is the integrity
counterpart of the companion interoperability design law \citep{baweja2026interop}
(standardize the ephemeris, not just the format), and it rests on the same open,
fully reproducible substrate as the companion datum-identifiability paper
\citep{baweja2026datum}. The proven and cited content carries over unchanged to any real
deployment; before any absolute CMPL number is trusted, a real multi-provider
constellation must substitute flight data for the representative ephemeris analog used
here.


\section*{Data and code availability}
\label{sec:availability}
The reference substrate for this paper is part of the open-source \code{kshana} engine
\citep{baweja2026kshana} (\url{https://github.com/ashfordeOU/kshana}; AGPL-3.0-only). The
geometry-matrix construction, the range/parity projectors, the parity/least-squares
split (\code{blind\_dx}, \code{blind\_fraction}, \code{detectable\_norm}), the
common-mode covariance, and the CMPL/envelope assembly are in the module
\code{lunar\_common\_mode}. Every number in this paper is regenerated by
\begin{quote}
  \code{cargo run -{}-example p3\_common\_mode\_integrity}
\end{quote}
which prints Part~A (the analytic-identity check, Eq.~\ref{eq:partA}), Part~B (the real
DE440/INPOP21a/EPM2021 split of Table~\ref{tab:anchor}), and Part~C (the modelled CMPL
integrity envelope of Table~\ref{tab:envelope}), directly from the committed engine and
the vendored fixtures. The independent NumPy oracle is
\code{scripts/gen\_common\_mode\_ref.py}, which reads the DE440/INPOP21a/EPM2021 kernels
through IMCCE \code{calceph}, forms the per-epoch geocentric-Moon common shift, and
emits the reference split by an independent linear-algebra path; its output and the
sampled positions are vendored at \code{tests/fixtures/lunar\_common\_mode/}
(\code{moon\_geo.csv}, \code{reference.json}, \code{NOTICE.md}). The ephemeris kernels
themselves are not vendored (they are large and their redistribution terms vary); the
notice records their public download URLs. The \Validated{} row---the engine's split
matching the NumPy oracle to relative and absolute error below $10^{-3}$---is asserted in
the continuous-integration suite by
\code{tests/lunar\_common\_mode\_integrity\_reference.rs}, which fails the build on any
mismatch. All frame, timescale, and sampling conventions are recorded in
Appendix~\ref{app:methods} and the vendored notice. A productized operational
ARAIM$+$CMPL integrity monitor---as opposed to the open reference substrate that
reproduces the paper's numbers---is a commercial capability of \code{kshana-pro} and is
not part of this open release.


\section*{Acknowledgements}
The numerical substrate underlying this work is part of the open-source
\code{kshana} engine. The lunar ephemerides differenced here are the property of
their respective institutions (JPL, IMCCE, IAA~RAS), gratefully acknowledged. This
research received no specific grant from any funding agency in the public,
commercial, or not-for-profit sectors.

\bibliographystyle{plainnat}
\bibliography{refs}

\appendix
\section{Methods and modelled assumptions}
\label{app:methods}

This appendix records the measurement-model conventions, the ephemeris handling, the
constellation model, and the numerical method behind the results. All are consistent
with the vendored provenance notice (Section~\ref{sec:availability}).

\subsection{Measurement model and projectors}
\label{app:model}

The geometry matrix $\Gmat$ of \eqref{eq:measmodel} uses the sign convention that the
position partial of the $i$-th pseudorange residual is the negative user-to-satellite
LOS unit vector $-\los_i$ and the clock partial is unity, so that a satellite-position
error $\ds_i$ enters the residual as $+\los_i\cdot\ds_i$ (an ephemeris error appears as
a range error along the sightline). The range and parity projectors
\eqref{eq:projectors} are formed from the thin QR factorization of $\Gmat$: with
$\Gmat=\bm Q\bm R$, $\bm Q\in\R^{n\times4}$ orthonormal, $\Ppar=\bm Q\bm Q^\top$ and
$\Pperp=\bm I_n-\bm Q\bm Q^\top$, which is numerically preferable to forming
$(\Gmat^\top\Gmat)^{-1}$ explicitly and is the route both the engine and the NumPy oracle
take. A parity matrix $\Smat$ (any orthonormal basis of $\rng(\Gmat)^\perp$, e.g.\ from
the full QR) gives the same $\Pperp=\Smat^\top\Smat$; the split quantities
\eqref{eq:blinddx}--\eqref{eq:detnorm} are basis-independent and we report them as such.

\subsection{Ephemeris provenance and access}
\label{app:provenance}

Three independent authoritative ephemeris families supply the real common-shift
magnitude, read through the IMCCE \code{calceph} library so that all three are accessed
by one interface in a single frame and timescale:
\begin{itemize}
  \item \textbf{DE440} (JPL/NASA), kernel \code{de440s.bsp} \citep{park2021de440}, from
  \url{https://naif.jpl.nasa.gov/pub/naif/generic_kernels/spk/planets/de440s.bsp};
  \item \textbf{INPOP21a} (IMCCE / Observatoire de Paris)
  \citep{fienga2021inpop21a}, from \url{https://ftp.imcce.fr/pub/ephem/planets/inpop21a/};
  \item \textbf{EPM2021} (IAA~RAS), kernel \code{epm2021.bsp}
  \citep{pitjeva2021epm2021}, from \url{https://ftp.iaaras.ru/pub/epm/EPM2021/SPICE/}.
\end{itemize}
Following the vendoring convention of the companion papers \citep{baweja2026datum,%
baweja2026interop}, the kernels are \emph{not} redistributed; only the sampled derived
positions and the computed reference split are vendored. Published scientific ephemeris
positions are factual constants and are not copyrightable; the derived sample points are
cited with full attribution to JPL, IMCCE, and IAA~RAS.

\subsection{Frame, timescale, and sampling}
\label{app:sampling}

The geocentric Moon position (NAIF body $301$ with respect to $399$) is sampled from
each family over 2024-01-01 to 2025-12-31 in Barycentric Dynamical Time (TDB) at two-day
cadence ($366$ epochs), in the ICRF/J2000 frame, so the three families are directly
comparable without a frame or timescale conversion. For each provider pair and epoch, the
per-epoch position difference is the common frame/ephemeris shift $\ds$ applied to the
constellation. Positions are rounded to six decimals (micrometre resolution) before the
fit in both the engine and the oracle, so the comparison tests the two \emph{algorithms}
at byte-identical inputs and is not sensitive to sampling-level differences. The record
count is $732$ ($2$ provider pairs $\times366$ epochs); Table~\ref{tab:anchor} reports
the per-pair medians.

\subsection{The modelled constellation}
\label{app:constellation}

The geometry matrix $\Gmat$ is built from a modelled lunar navigation constellation
(satellite LOS unit vectors to a representative user), used to project the common shift
$\ds$ into measurement space. This constellation is \Modelled{}: it stands in for a real
provider's beacon geometry, which has not flown. The blind-subspace membership of the
common shift (Lemma~\ref{lem:rigid}) is \emph{independent} of the constellation geometry
and of the user position---it holds for any $\Gmat$ of the form \eqref{eq:measmodel}---so
the modelled geometry affects only the partition of the fixed $\approx\!2\,$m magnitude
between the horizontal and vertical position components, not the fact of blindness. The
horizontal/vertical split feeding $\sigma_H$ in \eqref{eq:sigmaH} is modelled as
isotropic ($\Sig_{\mathrm{cm}}=\tfrac{1}{3}\sigma_{\mathrm{cm}}^2\bm I_3$); a specific
constellation and user geometry would replace this isotropic assumption with the actual
position-domain covariance, shifting $\sigma_H$ but not the mechanism.

\subsection{The CMPL and envelope}
\label{app:cmpl}

The common-mode one-sigma $\sigma_{\mathrm{cm}}=\SI{2}{m}$ is fixed to the real
inter-ephemeris floor of Section~\ref{sec:anchor} and \citep{baweja2026interop}. The
horizontal one-sigma follows from the isotropic assumption as
$\sigma_H=\sqrt{2/3}\,\sigma_{\mathrm{cm}}$, and the CMPL is the $2$-D horizontal Rayleigh
inflation $K_H\sigma_H$ with the representative multiplier $K_H=6.0$, the value used at
civil-aviation horizontal integrity risks in ARAIM \citep{blanch2015,wgc2016araim}. The
representative single-provider $\mathrm{HPL}_{\mathrm{ARAIM}}=\SI{15}{m}$ is an
illustrative placeholder, not a computed lunar protection level; it enters only to
exhibit the envelope's structure. The total is the triangle-inequality (arithmetic) sum
$\mathrm{HPL}_{\mathrm{ARAIM}}+\mathrm{CMPL}$ rather than a root-sum-square, because the
single-provider and cross-provider position-domain errors are not guaranteed independent
and can align (Section~\ref{sec:envelope}); the arithmetic sum is the conservative bound
and the root-sum-square would overclaim integrity. The envelope is valid only where the
cross-provider common-mode is not already absorbed into a provider's
$\sigma_{\mathrm{URA}}$ overbound (Section~\ref{sec:conditionality}); this
anti-double-counting condition is part of the model's stated scope.


\end{document}
